(i) Existence for $W$: $w_{i}f = 0$ since $w_{i}$ anticommutes with the other $w_{j}$ and $w_{i}^{2} = 0$. Uniqueness: on the grade-$k$ state $v_{i_{1}}\cdots v_{i_{k}}f$, $w_{j}$ acts as zero if $j \notin \{i_{1}, \ldots, i_{k}\}$ and otherwise, using $w_{j}v_{i} = \delta_{ij} - v_{i}w_{j}$ (\claimref{std-algebra.witt-basis-car}), removes the factor $v_{j}$ up to sign; so the common kernel of $w_{1}, \ldots, w_{m}$ on $S$ is the grade-zero line $\C f$. For $W'$ the same computation with the roles exchanged gives the top state, which every $v_{i}$ annihilates since $v_{i}^{2} = 0$. For a general $U$, Witt's theorem (\claimref{std-algebra.witt-extension}; as in Proposition \claimref{dark-matter.flag-insensitive-residue-no-flag-level}) supplies an isometry $\rho$ with $\rho(W) = U$, and $\rho$ is the twisted adjoint action of a versor $R$ (Lemma \claimref{primer.twisted-adjoint-action}), $\rho(x) = \pm RxR^{-1}$; then $U(Rf) = \pm R(Wf)R^{-1}R = 0$, so $Rf$ is a pure spinor of $U$, and it is the only one up to scale because $R$ is invertible on $S$. Chirality: $I$ commutes with every even element, so $I$ acts on each grade as a scalar; explicitly
\[
I \;\propto\; \prod_{i}(v_{i}w_{i} - w_{i}v_{i}) \;=\; \prod_{i}(2v_{i}w_{i} - 1),
\]
and the $i$-th factor is $-1$ on states without $v_{i}$ and $+1$ on states with it, so the grade-$k$ states have chirality $\epsilon(-1)^{m-k}$ with $\epsilon$ the normalisation; every pure spinor, being a versor image of $f$, is a chirality eigenstate by (ii).

(ii) In even dimension $e_{j}I = -Ie_{j}$ for every basis vector (\claimref{std-algebra.pseudoscalar-commutation}), so $I$ commutes with even versors and anticommutes with odd ones; if $I\psi = \chi\psi$ then $I(R\psi) = \pm R I\psi = \pm\chi(R\psi)$ with the sign that of $R$'s parity.

(iii) Both chiralities occur among pure spinors (reflect $W$ in any non-null vector), chirality is constant on each orbit of the connected group by (ii), and the full isometry group, which is the connected group together with one reflection, is transitive by Witt; so there are exactly two orbits, one per chirality. For the Witt pair, $f$ has grade $0$ and the top state grade $m$, so their chiralities are $\epsilon(-1)^{m}$ and $\epsilon$: opposite exactly when $m$ is odd. Odd dimension: $-1$ is an improper isometry ($\det(-1) = (-1)^{2m+1}$) that fixes every subspace, so every orbit of the full group is an orbit of the connected group, and Witt's transitivity gives one family; and on either half $S_{\pm}$ the extra generator $\ell$ acts as the Fock grade parity up to a scalar (script-checked in dimensions five and seven), while the spin-algebra bivector $e_{1}\ell$ anticommutes with $\ell$, so the spin algebra mixes the Fock grades within each irreducible module.

The script builds $I$, $f$ and the top state on the split real forms of dimensions four, six and eight in \texttt{GA.py} and checks that both are eigenstates of $I$ with eigenvalues in the ratio $(-1)^{m}$ exactly (only the ratio is invariant, the sign of $I$ being a convention; the chiralities agree at $m = 2, 4$ and differ at $m = 3$), that $I$ anticommutes with every generator in even dimension and commutes with every generator in dimensions five and seven, and exhibits the exchanging versors with their parities. The classical result, with the general criterion that two maximal isotropic subspaces lie in one family exactly when their intersection has dimension congruent to $m$ modulo two, is Chevalley's \cite[Ch.~III, \S1]{Chevalley1954}, see also \cite{LawsonMichelsohn1989}; only the Witt-pair case is used here, and it is proved here.
